\begingroup
\clearpage\section{Historias, escenarios y aceptación}\label{sec4}
\setlength{\parskip}{3pt}
\fontsize{10}{13}\selectfont
\renewcommand{\arraystretch}{1.0}
\titlespacing*{\subsection}{0pt}{8pt}{4pt}
Cada historia recoge dos escenarios seleccionados y una variante relevante. V indica válido e I inválido. Las reglas comunes y los datos de prueba se definen en el apartado 3.\par
\subsection*{HU-01. Registrar una cuenta}
\textbf{Requisito} R1\par
\textbf{Historia} Como persona no registrada quiero crear una cuenta para utilizar los servicios de movilidad.\par
\begingroup\fontsize{8.5}{10.5}\selectfont\setlength{\tabcolsep}{4pt}
\begin{tabularx}{\linewidth}{@{}p{1.35cm}YY@{}}\toprule
\textbf{Escenario} & \textbf{Entrada o condición} & \textbf{Resultado}\\\midrule
E-01-V & Datos válidos, correo nuevo & Cuenta creada\\[2pt]
E-01-I1 & Correo ya registrado & Alta rechazada\\[2pt]
E-01-I2 & Contraseña de 9 caracteres & Alta rechazada\\[2pt]
\bottomrule\end{tabularx}\endgroup\par
\begin{minipage}{\linewidth}
\begingroup\fontsize{9}{10.8}\selectfont\begin{tabularx}{\linewidth}{@{}p{2.2cm}Y@{}}\arrayrulecolor{black}\toprule
\textbf{E-01-V} & \textbf{Escenario válido \textperiodcentered{} Prueba PA-01-V}\\[1pt]\midrule
\textbf{Justificación} & comprueba el alta con datos válidos y su guardado en la base de datos.\\[1pt]\arrayrulecolor{black!25}\midrule[0.25pt]
\textbf{Análisis} & Probar los tamaños mínimos y máximos de los campos.\\[1pt]\arrayrulecolor{black!25}\midrule[0.25pt]
\textbf{Given} & no existe ana@example.com y el almacenamiento está disponible\\[1pt]\arrayrulecolor{black!25}\midrule[0.25pt]
\textbf{When} & se registra Ana con ana@example.com y la contraseña RutaSegura2026\\[1pt]\arrayrulecolor{black!25}\midrule[0.25pt]
\textbf{Then} & se crea una única cuenta y se devuelve su identificador, sin devolver la contraseña\\[1pt]\arrayrulecolor{black!25}\midrule[0.25pt]
\end{tabularx}\endgroup\par
\end{minipage}\par\vspace{4pt}
\begin{minipage}{\linewidth}
\begingroup\fontsize{9}{10.8}\selectfont\begin{tabularx}{\linewidth}{@{}p{2.2cm}Y@{}}\arrayrulecolor{black}\toprule
\textbf{E-01-I1} & \textbf{Escenario inválido \textperiodcentered{} Prueba PA-01-I}\\[1pt]\midrule
\textbf{Justificación} & un correo ya registrado impide crear otra cuenta con la misma identidad.\\[1pt]\arrayrulecolor{black!25}\midrule[0.25pt]
\textbf{Análisis} & Probar el mismo correo con mayúsculas y espacios exteriores.\\[1pt]\arrayrulecolor{black!25}\midrule[0.25pt]
\textbf{Given} & existe ana@example.com\\[1pt]\arrayrulecolor{black!25}\midrule[0.25pt]
\textbf{When} & se intenta registrar de nuevo ese correo con RutaSegura2026\\[1pt]\arrayrulecolor{black!25}\midrule[0.25pt]
\textbf{Then} & se rechaza por correo duplicado y el número de cuentas no cambia\\[1pt]\arrayrulecolor{black!25}\midrule[0.25pt]
\end{tabularx}\endgroup\par
\end{minipage}\par\vspace{4pt}
\subsection*{HU-02. Iniciar sesión}
\textbf{Requisito} R1\par
\textbf{Historia} Como persona registrada quiero iniciar sesión para acceder a mis datos.\par
\begingroup\fontsize{8.5}{10.5}\selectfont\setlength{\tabcolsep}{4pt}
\begin{tabularx}{\linewidth}{@{}p{1.35cm}YY@{}}\toprule
\textbf{Escenario} & \textbf{Entrada o condición} & \textbf{Resultado}\\\midrule
E-02-V & Credenciales correctas & Sesión iniciada\\[2pt]
E-02-I1 & Contraseña incorrecta & Acceso rechazado\\[2pt]
E-02-I2 & Campos vacíos & Acceso rechazado\\[2pt]
\bottomrule\end{tabularx}\endgroup\par
\begin{minipage}{\linewidth}
\begingroup\fontsize{9}{10.8}\selectfont\begin{tabularx}{\linewidth}{@{}p{2.2cm}Y@{}}\arrayrulecolor{black}\toprule
\textbf{E-02-V} & \textbf{Escenario válido \textperiodcentered{} Prueba PA-02-V}\\[1pt]\midrule
\textbf{Justificación} & comprueba que el acceso corresponde a la cuenta indicada.\\[1pt]\arrayrulecolor{black!25}\midrule[0.25pt]
\textbf{Análisis} & Comprobar correo y contraseña correctos e incorrectos.\\[1pt]\arrayrulecolor{black!25}\midrule[0.25pt]
\textbf{Given} & existe Ana con ana@example.com y RutaSegura2026\\[1pt]\arrayrulecolor{black!25}\midrule[0.25pt]
\textbf{When} & inicia sesión con esas credenciales\\[1pt]\arrayrulecolor{black!25}\midrule[0.25pt]
\textbf{Then} & se crea una sesión vinculada a Ana y se recuperan únicamente sus datos\\[1pt]\arrayrulecolor{black!25}\midrule[0.25pt]
\end{tabularx}\endgroup\par
\end{minipage}\par\vspace{4pt}
\begin{minipage}{\linewidth}
\begingroup\fontsize{9}{10.8}\selectfont\begin{tabularx}{\linewidth}{@{}p{2.2cm}Y@{}}\arrayrulecolor{black}\toprule
\textbf{E-02-I1} & \textbf{Escenario inválido \textperiodcentered{} Prueba PA-02-I}\\[1pt]\midrule
\textbf{Justificación} & una contraseña errónea no debe conceder acceso ni revelar datos.\\[1pt]\arrayrulecolor{black!25}\midrule[0.25pt]
\textbf{Análisis} & Repetir con correo inexistente y campos vacíos.\\[1pt]\arrayrulecolor{black!25}\midrule[0.25pt]
\textbf{Given} & existe la cuenta ana@example.com con contraseña RutaSegura2026\\[1pt]\arrayrulecolor{black!25}\midrule[0.25pt]
\textbf{When} & se inicia sesión con ana@example.com y OtraClave2026\\[1pt]\arrayrulecolor{black!25}\midrule[0.25pt]
\textbf{Then} & se informa de credenciales incorrectas y no se crea una sesión\\[1pt]\arrayrulecolor{black!25}\midrule[0.25pt]
\end{tabularx}\endgroup\par
\end{minipage}\par\vspace{4pt}
\clearpage
\subsection*{HU-03. Cerrar sesión}
\textbf{Requisito} R1\par
\textbf{Historia} Como usuario autenticado quiero cerrar mi sesión para impedir el acceso posterior desde ella.\par
\begingroup\fontsize{8.5}{10.5}\selectfont\setlength{\tabcolsep}{4pt}
\begin{tabularx}{\linewidth}{@{}p{1.35cm}YY@{}}\toprule
\textbf{Escenario} & \textbf{Entrada o condición} & \textbf{Resultado}\\\midrule
E-03-V & Sesión activa & Sesión cerrada, datos conservados\\[2pt]
E-03-I1 & Sesión ya cerrada & Cierre rechazado\\[2pt]
E-03-V2 & Volver a acceder con credenciales correctas & Nueva sesión\\[2pt]
\bottomrule\end{tabularx}\endgroup\par
\begin{minipage}{\linewidth}
\begingroup\fontsize{9}{10.8}\selectfont\begin{tabularx}{\linewidth}{@{}p{2.2cm}Y@{}}\arrayrulecolor{black}\toprule
\textbf{E-03-V} & \textbf{Escenario válido \textperiodcentered{} Prueba PA-03-V}\\[1pt]\midrule
\textbf{Justificación} & comprueba que la sesión cerrada deja de dar acceso a los datos.\\[1pt]\arrayrulecolor{black!25}\midrule[0.25pt]
\textbf{Análisis} & Comprobar otra sesión abierta de la misma cuenta.\\[1pt]\arrayrulecolor{black!25}\midrule[0.25pt]
\textbf{Given} & Ana tiene una sesión activa S1\\[1pt]\arrayrulecolor{black!25}\midrule[0.25pt]
\textbf{When} & cierra S1 y vuelve a solicitar sus datos con ese código de sesión\\[1pt]\arrayrulecolor{black!25}\midrule[0.25pt]
\textbf{Then} & S1 deja de dar acceso y los datos de Ana se conservan\\[1pt]\arrayrulecolor{black!25}\midrule[0.25pt]
\end{tabularx}\endgroup\par
\end{minipage}\par\vspace{4pt}
\begin{minipage}{\linewidth}
\begingroup\fontsize{9}{10.8}\selectfont\begin{tabularx}{\linewidth}{@{}p{2.2cm}Y@{}}\arrayrulecolor{black}\toprule
\textbf{E-03-I1} & \textbf{Escenario inválido \textperiodcentered{} Prueba PA-03-I}\\[1pt]\midrule
\textbf{Justificación} & no se puede cerrar una sesión que ya está cerrada.\\[1pt]\arrayrulecolor{black!25}\midrule[0.25pt]
\textbf{Análisis} & Comprobar que los datos guardados se conservan.\\[1pt]\arrayrulecolor{black!25}\midrule[0.25pt]
\textbf{Given} & Ana ha cerrado su sesión S1\\[1pt]\arrayrulecolor{black!25}\midrule[0.25pt]
\textbf{When} & intenta cerrar S1 de nuevo\\[1pt]\arrayrulecolor{black!25}\midrule[0.25pt]
\textbf{Then} & se informa de que la sesión no está activa y sus datos se conservan\\[1pt]\arrayrulecolor{black!25}\midrule[0.25pt]
\end{tabularx}\endgroup\par
\end{minipage}\par\vspace{4pt}
\subsection*{HU-04. Eliminar la cuenta}
\textbf{Requisito} R1\par
\textbf{Historia} Como usuario quiero eliminar mi cuenta para borrar sus datos asociados.\par
\begingroup\fontsize{8.5}{10.5}\selectfont\setlength{\tabcolsep}{4pt}
\begin{tabularx}{\linewidth}{@{}p{1.35cm}YY@{}}\toprule
\textbf{Escenario} & \textbf{Entrada o condición} & \textbf{Resultado}\\\midrule
E-04-V & Contraseña y confirmación correctas & Cuenta y datos eliminados\\[2pt]
E-04-I1 & Contraseña incorrecta & Datos conservados\\[2pt]
E-04-V2 & Cancelar la eliminación & Datos conservados\\[2pt]
\bottomrule\end{tabularx}\endgroup\par
\begin{minipage}{\linewidth}
\begingroup\fontsize{9}{10.8}\selectfont\begin{tabularx}{\linewidth}{@{}p{2.2cm}Y@{}}\arrayrulecolor{black}\toprule
\textbf{E-04-V} & \textbf{Escenario válido \textperiodcentered{} Prueba PA-04-V}\\[1pt]\midrule
\textbf{Justificación} & comprueba que se eliminan la cuenta y todos sus datos asociados.\\[1pt]\arrayrulecolor{black!25}\midrule[0.25pt]
\textbf{Análisis} & Repetir con cuenta vacía y con favoritos guardados.\\[1pt]\arrayrulecolor{black!25}\midrule[0.25pt]
\textbf{Given} & Ana tiene la contraseña RutaSegura2026, dos lugares, un vehículo y una ruta guardada\\[1pt]\arrayrulecolor{black!25}\midrule[0.25pt]
\textbf{When} & confirma la eliminación con RutaSegura2026\\[1pt]\arrayrulecolor{black!25}\midrule[0.25pt]
\textbf{Then} & se eliminan cuenta, sesiones, lugares, vehículo, ruta y preferencias de Ana\\[1pt]\arrayrulecolor{black!25}\midrule[0.25pt]
\end{tabularx}\endgroup\par
\end{minipage}\par\vspace{4pt}
\begin{minipage}{\linewidth}
\begingroup\fontsize{9}{10.8}\selectfont\begin{tabularx}{\linewidth}{@{}p{2.2cm}Y@{}}\arrayrulecolor{black}\toprule
\textbf{E-04-I1} & \textbf{Escenario inválido \textperiodcentered{} Prueba PA-04-I}\\[1pt]\midrule
\textbf{Justificación} & impide un borrado irreversible sin comprobar la contraseña.\\[1pt]\arrayrulecolor{black!25}\midrule[0.25pt]
\textbf{Análisis} & Cancelar y comprobar también que la sesión sigue activa.\\[1pt]\arrayrulecolor{black!25}\midrule[0.25pt]
\textbf{Given} & Ana tiene la contraseña RutaSegura2026, dos lugares, un vehículo y una ruta guardada\\[1pt]\arrayrulecolor{black!25}\midrule[0.25pt]
\textbf{When} & confirma la eliminación con OtraClave2026\\[1pt]\arrayrulecolor{black!25}\midrule[0.25pt]
\textbf{Then} & se rechaza la operación y permanecen la cuenta y todos sus datos\\[1pt]\arrayrulecolor{black!25}\midrule[0.25pt]
\end{tabularx}\endgroup\par
\end{minipage}\par\vspace{4pt}
\clearpage
\subsection*{HU-05. Guardar por coordenadas}
\textbf{Requisito} R2\par
\textbf{Historia} Como usuario quiero dar de alta un lugar mediante coordenadas para usarlo en una ruta.\par
\begingroup\fontsize{8.5}{10.5}\selectfont\setlength{\tabcolsep}{4pt}
\begin{tabularx}{\linewidth}{@{}p{1.35cm}YY@{}}\toprule
\textbf{Escenario} & \textbf{Entrada o condición} & \textbf{Resultado}\\\midrule
E-05-V & Nombre y coordenadas válidos & Lugar guardado\\[2pt]
E-05-I1 & Latitud 91 & Lugar no guardado\\[2pt]
E-05-I2 & Longitud 181 & Lugar no guardado\\[2pt]
\bottomrule\end{tabularx}\endgroup\par
\begin{minipage}{\linewidth}
\begingroup\fontsize{9}{10.8}\selectfont\begin{tabularx}{\linewidth}{@{}p{2.2cm}Y@{}}\arrayrulecolor{black}\toprule
\textbf{E-05-V} & \textbf{Escenario válido \textperiodcentered{} Prueba PA-05-V}\\[1pt]\midrule
\textbf{Justificación} & comprueba el alta directa mediante coordenadas.\\[1pt]\arrayrulecolor{black!25}\midrule[0.25pt]
\textbf{Análisis} & Incluir los extremos de latitud y longitud admitidos.\\[1pt]\arrayrulecolor{black!25}\midrule[0.25pt]
\textbf{Given} & Ana tiene cero lugares\\[1pt]\arrayrulecolor{black!25}\midrule[0.25pt]
\textbf{When} & guarda Centro con latitud 39.9864 y longitud -0.0513\\[1pt]\arrayrulecolor{black!25}\midrule[0.25pt]
\textbf{Then} & aparece un lugar de Ana con ese nombre y esas coordenadas\\[1pt]\arrayrulecolor{black!25}\midrule[0.25pt]
\end{tabularx}\endgroup\par
\end{minipage}\par\vspace{4pt}
\begin{minipage}{\linewidth}
\begingroup\fontsize{9}{10.8}\selectfont\begin{tabularx}{\linewidth}{@{}p{2.2cm}Y@{}}\arrayrulecolor{black}\toprule
\textbf{E-05-I1} & \textbf{Escenario inválido \textperiodcentered{} Prueba PA-05-I}\\[1pt]\midrule
\textbf{Justificación} & evita guardar coordenadas fuera de los límites admitidos.\\[1pt]\arrayrulecolor{black!25}\midrule[0.25pt]
\textbf{Análisis} & Probar longitud 181 y valores que no sean números.\\[1pt]\arrayrulecolor{black!25}\midrule[0.25pt]
\textbf{Given} & Ana tiene cero lugares\\[1pt]\arrayrulecolor{black!25}\midrule[0.25pt]
\textbf{When} & guarda Centro con latitud 91 y longitud -0.0513\\[1pt]\arrayrulecolor{black!25}\midrule[0.25pt]
\textbf{Then} & se rechaza la latitud y la lista sigue vacía\\[1pt]\arrayrulecolor{black!25}\midrule[0.25pt]
\end{tabularx}\endgroup\par
\end{minipage}\par\vspace{4pt}
\subsection*{HU-06. Guardar por topónimo}
\textbf{Requisito} R2\par
\textbf{Historia} Como usuario quiero buscar un topónimo y guardar un resultado para identificar un lugar sin conocer sus coordenadas.\par
\begingroup\fontsize{8.5}{10.5}\selectfont\setlength{\tabcolsep}{4pt}
\begin{tabularx}{\linewidth}{@{}p{1.35cm}YY@{}}\toprule
\textbf{Escenario} & \textbf{Entrada o condición} & \textbf{Resultado}\\\midrule
E-06-V & Búsqueda con candidato seleccionado & Candidato guardado\\[2pt]
E-06-I1 & Topónimo de tres espacios & Búsqueda rechazada\\[2pt]
E-06-V2 & Búsqueda sin coincidencias & Lista vacía, sin alta\\[2pt]
\bottomrule\end{tabularx}\endgroup\par
\begin{minipage}{\linewidth}
\begingroup\fontsize{9}{10.8}\selectfont\begin{tabularx}{\linewidth}{@{}p{2.2cm}Y@{}}\arrayrulecolor{black}\toprule
\textbf{E-06-V} & \textbf{Escenario válido \textperiodcentered{} Prueba PA-06-V}\\[1pt]\midrule
\textbf{Justificación} & comprueba la búsqueda y el guardado del candidato elegido, aunque el servicio cambie el orden de resultados.\\[1pt]\arrayrulecolor{black!25}\midrule[0.25pt]
\textbf{Análisis} & Elegir entre varios resultados y comprobar sus coordenadas.\\[1pt]\arrayrulecolor{black!25}\midrule[0.25pt]
\textbf{Given} & Ana está autenticada y el servicio real de búsqueda de lugares está disponible\\[1pt]\arrayrulecolor{black!25}\midrule[0.25pt]
\textbf{When} & busca Universitat Jaume I, Castellón, selecciona un resultado y lo guarda como UJI\\[1pt]\arrayrulecolor{black!25}\midrule[0.25pt]
\textbf{Then} & se obtiene al menos un candidato con coordenadas válidas y se guarda el nombre UJI con las coordenadas del resultado seleccionado\\[1pt]\arrayrulecolor{black!25}\midrule[0.25pt]
\end{tabularx}\endgroup\par
\end{minipage}\par\vspace{4pt}
\begin{minipage}{\linewidth}
\begingroup\fontsize{9}{10.8}\selectfont\begin{tabularx}{\linewidth}{@{}p{2.2cm}Y@{}}\arrayrulecolor{black}\toprule
\textbf{E-06-I1} & \textbf{Escenario inválido \textperiodcentered{} Prueba PA-06-I}\\[1pt]\midrule
\textbf{Justificación} & una entrada vacía carece de significado y no debe consumir peticiones.\\[1pt]\arrayrulecolor{black!25}\midrule[0.25pt]
\textbf{Análisis} & Repetir con búsqueda vacía, sin resultados y de 201 caracteres.\\[1pt]\arrayrulecolor{black!25}\midrule[0.25pt]
\textbf{Given} & Ana tiene cero lugares\\[1pt]\arrayrulecolor{black!25}\midrule[0.25pt]
\textbf{When} & solicita buscar un topónimo formado por tres espacios\\[1pt]\arrayrulecolor{black!25}\midrule[0.25pt]
\textbf{Then} & se rechaza la búsqueda vacía antes de consultar el servicio y no se guarda ningún lugar\\[1pt]\arrayrulecolor{black!25}\midrule[0.25pt]
\end{tabularx}\endgroup\par
\end{minipage}\par\vspace{4pt}
\clearpage
\subsection*{HU-07. Consultar lugares}
\textbf{Requisito} R2\par
\textbf{Historia} Como usuario quiero consultar mis lugares para elegir los que utilizaré.\par
\begingroup\fontsize{8.5}{10.5}\selectfont\setlength{\tabcolsep}{4pt}
\begin{tabularx}{\linewidth}{@{}p{1.35cm}YY@{}}\toprule
\textbf{Escenario} & \textbf{Entrada o condición} & \textbf{Resultado}\\\midrule
E-07-V & Sesión activa, Centro y UJI guardados & Lista de lugares\\[2pt]
E-07-I1 & Sin sesión válida & Solicitar acceso\\[2pt]
E-07-V2 & Ningún lugar guardado & Lista vacía\\[2pt]
\bottomrule\end{tabularx}\endgroup\par
\begin{minipage}{\linewidth}
\begingroup\fontsize{9}{10.8}\selectfont\begin{tabularx}{\linewidth}{@{}p{2.2cm}Y@{}}\arrayrulecolor{black}\toprule
\textbf{E-07-V} & \textbf{Escenario válido \textperiodcentered{} Prueba PA-07-V}\\[1pt]\midrule
\textbf{Justificación} & comprueba que se recuperan los lugares guardados.\\[1pt]\arrayrulecolor{black!25}\midrule[0.25pt]
\textbf{Análisis} & Repetir con lista vacía y con lugares favoritos.\\[1pt]\arrayrulecolor{black!25}\midrule[0.25pt]
\textbf{Given} & Ana tiene Centro y UJI\\[1pt]\arrayrulecolor{black!25}\midrule[0.25pt]
\textbf{When} & Ana consulta su lista\\[1pt]\arrayrulecolor{black!25}\midrule[0.25pt]
\textbf{Then} & se devuelven Centro y UJI con sus coordenadas\\[1pt]\arrayrulecolor{black!25}\midrule[0.25pt]
\end{tabularx}\endgroup\par
\end{minipage}\par\vspace{4pt}
\begin{minipage}{\linewidth}
\begingroup\fontsize{9}{10.8}\selectfont\begin{tabularx}{\linewidth}{@{}p{2.2cm}Y@{}}\arrayrulecolor{black}\toprule
\textbf{E-07-I1} & \textbf{Escenario inválido \textperiodcentered{} Prueba PA-07-I}\\[1pt]\midrule
\textbf{Justificación} & evita revelar ubicaciones sin autenticación.\\[1pt]\arrayrulecolor{black!25}\midrule[0.25pt]
\textbf{Análisis} & Repetir con sesión cerrada y con sesión caducada.\\[1pt]\arrayrulecolor{black!25}\midrule[0.25pt]
\textbf{Given} & no hay sesión válida\\[1pt]\arrayrulecolor{black!25}\midrule[0.25pt]
\textbf{When} & se solicita la lista de Ana\\[1pt]\arrayrulecolor{black!25}\midrule[0.25pt]
\textbf{Then} & se deniega el acceso y no se muestran lugares\\[1pt]\arrayrulecolor{black!25}\midrule[0.25pt]
\end{tabularx}\endgroup\par
\end{minipage}\par\vspace{4pt}
\subsection*{HU-08. Eliminar un lugar}
\textbf{Requisito} R2\par
\textbf{Historia} Como usuario quiero eliminar un lugar propio para mantener actualizada mi lista.\par
\begingroup\fontsize{8.5}{10.5}\selectfont\setlength{\tabcolsep}{4pt}
\begin{tabularx}{\linewidth}{@{}p{1.35cm}YY@{}}\toprule
\textbf{Escenario} & \textbf{Entrada o condición} & \textbf{Resultado}\\\midrule
E-08-V & Centro y UJI, confirmar borrar Centro & Solo queda UJI\\[2pt]
E-08-I1 & Centro eliminado antes de confirmar & Aviso, UJI conservado\\[2pt]
E-08-V2 & Cancelar el borrado & Lugar conservado\\[2pt]
\bottomrule\end{tabularx}\endgroup\par
\begin{minipage}{\linewidth}
\begingroup\fontsize{9}{10.8}\selectfont\begin{tabularx}{\linewidth}{@{}p{2.2cm}Y@{}}\arrayrulecolor{black}\toprule
\textbf{E-08-V} & \textbf{Escenario válido \textperiodcentered{} Prueba PA-08-V}\\[1pt]\midrule
\textbf{Justificación} & comprueba que la eliminación afecta solo al registro elegido.\\[1pt]\arrayrulecolor{black!25}\midrule[0.25pt]
\textbf{Análisis} & Comprobar la cancelación y las rutas históricas del lugar.\\[1pt]\arrayrulecolor{black!25}\midrule[0.25pt]
\textbf{Given} & Ana tiene Centro y UJI\\[1pt]\arrayrulecolor{black!25}\midrule[0.25pt]
\textbf{When} & confirma eliminar Centro\\[1pt]\arrayrulecolor{black!25}\midrule[0.25pt]
\textbf{Then} & solo queda UJI y Centro deja de poder elegirse en nuevas rutas\\[1pt]\arrayrulecolor{black!25}\midrule[0.25pt]
\end{tabularx}\endgroup\par
\end{minipage}\par\vspace{4pt}
\begin{minipage}{\linewidth}
\begingroup\fontsize{9}{10.8}\selectfont\begin{tabularx}{\linewidth}{@{}p{2.2cm}Y@{}}\arrayrulecolor{black}\toprule
\textbf{E-08-I1} & \textbf{Escenario inválido \textperiodcentered{} Prueba PA-08-I}\\[1pt]\midrule
\textbf{Justificación} & el lugar puede dejar de estar disponible antes de confirmar.\\[1pt]\arrayrulecolor{black!25}\midrule[0.25pt]
\textbf{Análisis} & Comprobar que los demás lugares se conservan.\\[1pt]\arrayrulecolor{black!25}\midrule[0.25pt]
\textbf{Given} & Ana tenía Centro y UJI, pero Centro se ha eliminado desde otra sesión mientras el diálogo seguía abierto\\[1pt]\arrayrulecolor{black!25}\midrule[0.25pt]
\textbf{When} & confirma eliminar Centro desde el diálogo anterior\\[1pt]\arrayrulecolor{black!25}\midrule[0.25pt]
\textbf{Then} & se informa de que Centro ya no está disponible y UJI permanece\\[1pt]\arrayrulecolor{black!25}\midrule[0.25pt]
\end{tabularx}\endgroup\par
\end{minipage}\par\vspace{4pt}
\clearpage
\subsection*{HU-09. Elegir lugares en el mapa}
\textbf{Requisito} R2\par
\textbf{Historia} Como usuario quiero seleccionar un punto en el mapa y guardarlo para fijar visualmente un origen o un destino.\par
\begingroup\fontsize{8.5}{10.5}\selectfont\setlength{\tabcolsep}{4pt}
\begin{tabularx}{\linewidth}{@{}p{1.35cm}YY@{}}\toprule
\textbf{Escenario} & \textbf{Entrada o condición} & \textbf{Resultado}\\\midrule
E-09-V & Punto seleccionado y nombre válido & Lugar guardado\\[2pt]
E-09-I1 & Sin punto seleccionado & Lugar no guardado\\[2pt]
E-09-V2 & Mover el marcador antes de guardar & Se guarda la última posición\\[2pt]
\bottomrule\end{tabularx}\endgroup\par
\begin{minipage}{\linewidth}
\begingroup\fontsize{9}{10.8}\selectfont\begin{tabularx}{\linewidth}{@{}p{2.2cm}Y@{}}\arrayrulecolor{black}\toprule
\textbf{E-09-V} & \textbf{Escenario válido \textperiodcentered{} Prueba PA-09-V}\\[1pt]\midrule
\textbf{Justificación} & comprueba la selección de un punto en el mapa exigida en el enunciado.\\[1pt]\arrayrulecolor{black!25}\midrule[0.25pt]
\textbf{Análisis} & Cambiar el punto antes de confirmar el alta.\\[1pt]\arrayrulecolor{black!25}\midrule[0.25pt]
\textbf{Given} & Ana visualiza el mapa y tiene cero lugares\\[1pt]\arrayrulecolor{black!25}\midrule[0.25pt]
\textbf{When} & selecciona el punto (39.9864,-0.0513), le da el nombre Centro y confirma el alta\\[1pt]\arrayrulecolor{black!25}\midrule[0.25pt]
\textbf{Then} & se guarda Centro en esas coordenadas y puede seleccionarse como extremo\\[1pt]\arrayrulecolor{black!25}\midrule[0.25pt]
\end{tabularx}\endgroup\par
\end{minipage}\par\vspace{4pt}
\begin{minipage}{\linewidth}
\begingroup\fontsize{9}{10.8}\selectfont\begin{tabularx}{\linewidth}{@{}p{2.2cm}Y@{}}\arrayrulecolor{black}\toprule
\textbf{E-09-I1} & \textbf{Escenario inválido \textperiodcentered{} Prueba PA-09-I}\\[1pt]\midrule
\textbf{Justificación} & impide guardar un lugar sin haber elegido sus coordenadas.\\[1pt]\arrayrulecolor{black!25}\midrule[0.25pt]
\textbf{Análisis} & Repetir cuando falta solo una de las coordenadas.\\[1pt]\arrayrulecolor{black!25}\midrule[0.25pt]
\textbf{Given} & Ana abre el alta con nombre Centro y sin coordenadas ni punto seleccionado\\[1pt]\arrayrulecolor{black!25}\midrule[0.25pt]
\textbf{When} & intenta guardar el lugar\\[1pt]\arrayrulecolor{black!25}\midrule[0.25pt]
\textbf{Then} & se solicita completar las coordenadas y no se crea ningún lugar\\[1pt]\arrayrulecolor{black!25}\midrule[0.25pt]
\end{tabularx}\endgroup\par
\end{minipage}\par\vspace{4pt}
\subsection*{HU-10. Dar de alta un vehículo}
\textbf{Requisito} R3\par
\textbf{Historia} Como usuario quiero registrar un vehículo para calcular rutas con su consumo.\par
\begingroup\fontsize{8.5}{10.5}\selectfont\setlength{\tabcolsep}{4pt}
\begin{tabularx}{\linewidth}{@{}p{1.35cm}YY@{}}\toprule
\textbf{Escenario} & \textbf{Entrada o condición} & \textbf{Resultado}\\\midrule
E-10-V & Gasolina, consumo 6 l/100 km & Vehículo guardado\\[2pt]
E-10-I1 & Consumo -1 & Alta rechazada\\[2pt]
E-10-V2 & Eléctrico, consumo 15 kWh/100 km & Vehículo guardado\\[2pt]
\bottomrule\end{tabularx}\endgroup\par
\begin{minipage}{\linewidth}
\begingroup\fontsize{9}{10.8}\selectfont\begin{tabularx}{\linewidth}{@{}p{2.2cm}Y@{}}\arrayrulecolor{black}\toprule
\textbf{E-10-V} & \textbf{Escenario válido \textperiodcentered{} Prueba PA-10-V}\\[1pt]\midrule
\textbf{Justificación} & identifica las unidades necesarias para el coste.\\[1pt]\arrayrulecolor{black!25}\midrule[0.25pt]
\textbf{Análisis} & Repetir con diésel y electricidad, comprobando sus unidades.\\[1pt]\arrayrulecolor{black!25}\midrule[0.25pt]
\textbf{Given} & Ana no tiene vehículos\\[1pt]\arrayrulecolor{black!25}\midrule[0.25pt]
\textbf{When} & registra Compacto, gasolina, consumo 6,0 l/100 km\\[1pt]\arrayrulecolor{black!25}\midrule[0.25pt]
\textbf{Then} & se guarda un vehículo propio con esos valores y queda disponible en el planificador\\[1pt]\arrayrulecolor{black!25}\midrule[0.25pt]
\end{tabularx}\endgroup\par
\end{minipage}\par\vspace{4pt}
\begin{minipage}{\linewidth}
\begingroup\fontsize{9}{10.8}\selectfont\begin{tabularx}{\linewidth}{@{}p{2.2cm}Y@{}}\arrayrulecolor{black}\toprule
\textbf{E-10-I1} & \textbf{Escenario inválido \textperiodcentered{} Prueba PA-10-I}\\[1pt]\midrule
\textbf{Justificación} & un consumo negativo produciría estimaciones sin sentido.\\[1pt]\arrayrulecolor{black!25}\midrule[0.25pt]
\textbf{Análisis} & Probar consumos de 0, de 0,09 y de 100,01.\\[1pt]\arrayrulecolor{black!25}\midrule[0.25pt]
\textbf{Given} & Ana no tiene vehículos\\[1pt]\arrayrulecolor{black!25}\midrule[0.25pt]
\textbf{When} & registra Compacto, gasolina, consumo -1 l/100 km\\[1pt]\arrayrulecolor{black!25}\midrule[0.25pt]
\textbf{Then} & se rechaza el consumo y la lista permanece vacía\\[1pt]\arrayrulecolor{black!25}\midrule[0.25pt]
\end{tabularx}\endgroup\par
\end{minipage}\par\vspace{4pt}
\clearpage
\subsection*{HU-11. Consultar vehículos}
\textbf{Requisito} R3\par
\textbf{Historia} Como usuario quiero consultar mis vehículos para seleccionar el apropiado.\par
\begingroup\fontsize{8.5}{10.5}\selectfont\setlength{\tabcolsep}{4pt}
\begin{tabularx}{\linewidth}{@{}p{1.35cm}YY@{}}\toprule
\textbf{Escenario} & \textbf{Entrada o condición} & \textbf{Resultado}\\\midrule
E-11-V & Sesión activa y vehículos guardados & Lista de vehículos\\[2pt]
E-11-I1 & Sesión caducada & Solicitar acceso\\[2pt]
E-11-V2 & Ningún vehículo guardado & Lista vacía\\[2pt]
\bottomrule\end{tabularx}\endgroup\par
\begin{minipage}{\linewidth}
\begingroup\fontsize{9}{10.8}\selectfont\begin{tabularx}{\linewidth}{@{}p{2.2cm}Y@{}}\arrayrulecolor{black}\toprule
\textbf{E-11-V} & \textbf{Escenario válido \textperiodcentered{} Prueba PA-11-V}\\[1pt]\midrule
\textbf{Justificación} & presenta los datos que permiten distinguir las opciones.\\[1pt]\arrayrulecolor{black!25}\midrule[0.25pt]
\textbf{Análisis} & Comprobar nombre, tipo y consumo de cada vehículo.\\[1pt]\arrayrulecolor{black!25}\midrule[0.25pt]
\textbf{Given} & Ana tiene Compacto y Eléctrico\\[1pt]\arrayrulecolor{black!25}\midrule[0.25pt]
\textbf{When} & Ana consulta sus vehículos\\[1pt]\arrayrulecolor{black!25}\midrule[0.25pt]
\textbf{Then} & se devuelven los dos vehículos con sus tipos y consumos\\[1pt]\arrayrulecolor{black!25}\midrule[0.25pt]
\end{tabularx}\endgroup\par
\end{minipage}\par\vspace{4pt}
\begin{minipage}{\linewidth}
\begingroup\fontsize{9}{10.8}\selectfont\begin{tabularx}{\linewidth}{@{}p{2.2cm}Y@{}}\arrayrulecolor{black}\toprule
\textbf{E-11-I1} & \textbf{Escenario inválido \textperiodcentered{} Prueba PA-11-I}\\[1pt]\midrule
\textbf{Justificación} & no expone información desde una sesión caducada.\\[1pt]\arrayrulecolor{black!25}\midrule[0.25pt]
\textbf{Análisis} & Repetir con una cuenta que ya haya cerrado sesión.\\[1pt]\arrayrulecolor{black!25}\midrule[0.25pt]
\textbf{Given} & la sesión de Ana ha caducado\\[1pt]\arrayrulecolor{black!25}\midrule[0.25pt]
\textbf{When} & consulta los vehículos\\[1pt]\arrayrulecolor{black!25}\midrule[0.25pt]
\textbf{Then} & se solicita iniciar sesión y no se devuelve la lista\\[1pt]\arrayrulecolor{black!25}\midrule[0.25pt]
\end{tabularx}\endgroup\par
\end{minipage}\par\vspace{4pt}
\subsection*{HU-12. Eliminar un vehículo}
\textbf{Requisito} R3\par
\textbf{Historia} Como usuario quiero eliminar un vehículo que ya no utilizo para evitar seleccionarlo.\par
\begingroup\fontsize{8.5}{10.5}\selectfont\setlength{\tabcolsep}{4pt}
\begin{tabularx}{\linewidth}{@{}p{1.35cm}YY@{}}\toprule
\textbf{Escenario} & \textbf{Entrada o condición} & \textbf{Resultado}\\\midrule
E-12-V & Eliminar Compacto predeterminado & Vehículo y referencia eliminados\\[2pt]
E-12-I1 & Compacto eliminado antes de confirmar & Aviso, resto conservado\\[2pt]
E-12-V2 & Cancelar el borrado & Vehículo conservado\\[2pt]
\bottomrule\end{tabularx}\endgroup\par
\begin{minipage}{\linewidth}
\begingroup\fontsize{9}{10.8}\selectfont\begin{tabularx}{\linewidth}{@{}p{2.2cm}Y@{}}\arrayrulecolor{black}\toprule
\textbf{E-12-V} & \textbf{Escenario válido \textperiodcentered{} Prueba PA-12-V}\\[1pt]\midrule
\textbf{Justificación} & comprueba también la integridad de las preferencias.\\[1pt]\arrayrulecolor{black!25}\midrule[0.25pt]
\textbf{Análisis} & Repetir con un vehículo favorito y uno sin uso.\\[1pt]\arrayrulecolor{black!25}\midrule[0.25pt]
\textbf{Given} & Compacto es el vehículo predeterminado de Ana\\[1pt]\arrayrulecolor{black!25}\midrule[0.25pt]
\textbf{When} & Ana confirma eliminar Compacto\\[1pt]\arrayrulecolor{black!25}\midrule[0.25pt]
\textbf{Then} & se borra el vehículo y la referencia predeterminada queda vacía\\[1pt]\arrayrulecolor{black!25}\midrule[0.25pt]
\end{tabularx}\endgroup\par
\end{minipage}\par\vspace{4pt}
\begin{minipage}{\linewidth}
\begingroup\fontsize{9}{10.8}\selectfont\begin{tabularx}{\linewidth}{@{}p{2.2cm}Y@{}}\arrayrulecolor{black}\toprule
\textbf{E-12-I1} & \textbf{Escenario inválido \textperiodcentered{} Prueba PA-12-I}\\[1pt]\midrule
\textbf{Justificación} & un vehículo ya eliminado no puede borrarse de nuevo.\\[1pt]\arrayrulecolor{black!25}\midrule[0.25pt]
\textbf{Análisis} & Comprobar que el resto de vehículos y preferencias se conserva.\\[1pt]\arrayrulecolor{black!25}\midrule[0.25pt]
\textbf{Given} & Ana tenía Compacto y Eléctrico, pero Compacto se ha eliminado desde otra sesión\\[1pt]\arrayrulecolor{black!25}\midrule[0.25pt]
\textbf{When} & confirma eliminar Compacto desde la lista anterior\\[1pt]\arrayrulecolor{black!25}\midrule[0.25pt]
\textbf{Then} & se informa de que Compacto ya no está disponible y Eléctrico permanece\\[1pt]\arrayrulecolor{black!25}\midrule[0.25pt]
\end{tabularx}\endgroup\par
\end{minipage}\par\vspace{4pt}
\clearpage
\subsection*{HU-13. Modificar un vehículo}
\textbf{Requisito} R3\par
\textbf{Historia} Como usuario quiero modificar el nombre, el tipo o el consumo de mi vehículo para mantener sus datos actualizados.\par
\begingroup\fontsize{8.5}{10.5}\selectfont\setlength{\tabcolsep}{4pt}
\begin{tabularx}{\linewidth}{@{}p{1.35cm}YY@{}}\toprule
\textbf{Escenario} & \textbf{Entrada o condición} & \textbf{Resultado}\\\midrule
E-13-V & Cambiar consumo a 5,5 l/100 km & Mismo vehículo actualizado\\[2pt]
E-13-I1 & Cambiar consumo a 0 & Consumo anterior conservado\\[2pt]
E-13-V2 & Cambiar de gasolina a eléctrico & Unidad kWh/100 km\\[2pt]
\bottomrule\end{tabularx}\endgroup\par
\begin{minipage}{\linewidth}
\begingroup\fontsize{9}{10.8}\selectfont\begin{tabularx}{\linewidth}{@{}p{2.2cm}Y@{}}\arrayrulecolor{black}\toprule
\textbf{E-13-V} & \textbf{Escenario válido \textperiodcentered{} Prueba PA-13-V}\\[1pt]\midrule
\textbf{Justificación} & conserva la identidad del vehículo al editar.\\[1pt]\arrayrulecolor{black!25}\midrule[0.25pt]
\textbf{Análisis} & Repetir cambiando el nombre y el tipo de vehículo.\\[1pt]\arrayrulecolor{black!25}\midrule[0.25pt]
\textbf{Given} & Compacto pertenece a Ana y consume 6,0 l/100 km\\[1pt]\arrayrulecolor{black!25}\midrule[0.25pt]
\textbf{When} & Ana cambia el consumo a 5,5 y guarda\\[1pt]\arrayrulecolor{black!25}\midrule[0.25pt]
\textbf{Then} & se mantiene el identificador y se registra 5,5 para los cálculos posteriores\\[1pt]\arrayrulecolor{black!25}\midrule[0.25pt]
\end{tabularx}\endgroup\par
\end{minipage}\par\vspace{4pt}
\begin{minipage}{\linewidth}
\begingroup\fontsize{9}{10.8}\selectfont\begin{tabularx}{\linewidth}{@{}p{2.2cm}Y@{}}\arrayrulecolor{black}\toprule
\textbf{E-13-I1} & \textbf{Escenario inválido \textperiodcentered{} Prueba PA-13-I}\\[1pt]\midrule
\textbf{Justificación} & comprueba que un consumo inválido no modifica el vehículo.\\[1pt]\arrayrulecolor{black!25}\midrule[0.25pt]
\textbf{Análisis} & Repetir con consumo superior a 100 y nombre vacío.\\[1pt]\arrayrulecolor{black!25}\midrule[0.25pt]
\textbf{Given} & Compacto pertenece a Ana y tiene un consumo de 6,0 l/100 km\\[1pt]\arrayrulecolor{black!25}\midrule[0.25pt]
\textbf{When} & Ana intenta guardar un consumo de 0 l/100 km\\[1pt]\arrayrulecolor{black!25}\midrule[0.25pt]
\textbf{Then} & se rechaza la edición y el consumo de Compacto permanece en 6,0 l/100 km\\[1pt]\arrayrulecolor{black!25}\midrule[0.25pt]
\end{tabularx}\endgroup\par
\end{minipage}\par\vspace{4pt}
\subsection*{HU-14. Calcular y visualizar una ruta}
\textbf{Requisito} R4\par
\textbf{Historia} Como usuario quiero obtener una ruta entre dos puntos elegidos en Explorar o lugares guardados para conocer su trayecto, distancia y duración.\par
\begingroup\fontsize{8.5}{10.5}\selectfont\setlength{\tabcolsep}{4pt}
\begin{tabularx}{\linewidth}{@{}p{1.35cm}YY@{}}\toprule
\textbf{Escenario} & \textbf{Entrada o condición} & \textbf{Resultado}\\\midrule
E-14-V & Extremos distintos y proveedor disponible & Trayecto y métricas\\[2pt]
E-14-I1 & Origen y destino coincidentes & Cálculo rechazado\\[2pt]
E-14-I2 & Proveedor no disponible & Error, sin ruta nueva\\[2pt]
\bottomrule\end{tabularx}\endgroup\par
\begin{minipage}{\linewidth}
\begingroup\fontsize{9}{10.8}\selectfont\begin{tabularx}{\linewidth}{@{}p{2.2cm}Y@{}}\arrayrulecolor{black}\toprule
\textbf{E-14-V} & \textbf{Escenario válido \textperiodcentered{} Prueba PA-14-V}\\[1pt]\midrule
\textbf{Justificación} & comprueba que el servicio devuelve un trayecto y unas medidas coherentes.\\[1pt]\arrayrulecolor{black!25}\midrule[0.25pt]
\textbf{Análisis} & Repetir en coche y bicicleta, con lugares guardados y puntos directos.\\[1pt]\arrayrulecolor{black!25}\midrule[0.25pt]
\textbf{Given} & Ana tiene Centro (39.9864,-0.0513) y UJI (39.9946,-0.0684), y el servicio real está disponible\\[1pt]\arrayrulecolor{black!25}\midrule[0.25pt]
\textbf{When} & calcula la ruta a pie entre Centro y UJI\\[1pt]\arrayrulecolor{black!25}\midrule[0.25pt]
\textbf{Then} & se devuelve una geometría con al menos dos puntos válidos, distancia positiva menor de 100 km y duración positiva\\[1pt]\arrayrulecolor{black!25}\midrule[0.25pt]
\end{tabularx}\endgroup\par
\end{minipage}\par\vspace{4pt}
\begin{minipage}{\linewidth}
\begingroup\fontsize{9}{10.8}\selectfont\begin{tabularx}{\linewidth}{@{}p{2.2cm}Y@{}}\arrayrulecolor{black}\toprule
\textbf{E-14-I1} & \textbf{Escenario inválido \textperiodcentered{} Prueba PA-14-I}\\[1pt]\midrule
\textbf{Justificación} & detecta una combinación de entradas incompatible con el trayecto definido.\\[1pt]\arrayrulecolor{black!25}\midrule[0.25pt]
\textbf{Análisis} & Probar coordenadas iguales aunque los nombres sean distintos.\\[1pt]\arrayrulecolor{black!25}\midrule[0.25pt]
\textbf{Given} & Ana tiene Centro\\[1pt]\arrayrulecolor{black!25}\midrule[0.25pt]
\textbf{When} & solicita una ruta de Centro a Centro\\[1pt]\arrayrulecolor{black!25}\midrule[0.25pt]
\textbf{Then} & se rechaza por extremos coincidentes y no se devuelve ni guarda una ruta nueva\\[1pt]\arrayrulecolor{black!25}\midrule[0.25pt]
\end{tabularx}\endgroup\par
\end{minipage}\par\vspace{4pt}
\clearpage
\subsection*{HU-15. Estimar coste de combustible}
\textbf{Requisito} R4\par
\textbf{Historia} Como usuario quiero conocer el coste de combustible para comparar el gasto de mi desplazamiento.\par
\begingroup\fontsize{8.5}{10.5}\selectfont\setlength{\tabcolsep}{4pt}
\begin{tabularx}{\linewidth}{@{}p{1.35cm}YY@{}}\toprule
\textbf{Escenario} & \textbf{Entrada o condición} & \textbf{Resultado}\\\midrule
E-15-V & Gasolina, consumo y precio válidos & Coste según distancia\\[2pt]
E-15-I1 & Precio manual -1 euros/l & Precio rechazado\\[2pt]
E-15-V2 & Precio público no disponible & Ruta consultable sin coste\\[2pt]
\bottomrule\end{tabularx}\endgroup\par
\begin{minipage}{\linewidth}
\begingroup\fontsize{9}{10.8}\selectfont\begin{tabularx}{\linewidth}{@{}p{2.2cm}Y@{}}\arrayrulecolor{black}\toprule
\textbf{E-15-V} & \textbf{Escenario válido \textperiodcentered{} Prueba PA-15-V}\\[1pt]\midrule
\textbf{Justificación} & usa la distancia real obtenida y un precio configurado para verificar el cálculo sin exigir una ruta de longitud fija.\\[1pt]\arrayrulecolor{black!25}\midrule[0.25pt]
\textbf{Análisis} & Repetir con diésel y revisar el redondeo final.\\[1pt]\arrayrulecolor{black!25}\midrule[0.25pt]
\textbf{Given} & Ana tiene Centro y UJI, Compacto con consumo de 6 l/100 km y precio manual de gasolina de 1,60 \texteuro{}/l\\[1pt]\arrayrulecolor{black!25}\midrule[0.25pt]
\textbf{When} & calcula una ruta en coche entre Centro y UJI con el servicio de rutas disponible\\[1pt]\arrayrulecolor{black!25}\midrule[0.25pt]
\textbf{Then} & el coste devuelto es (distancia en km / 100 \ensuremath{\times} 6 \ensuremath{\times} 1,60) \texteuro{} redondeado a céntimos, con el origen manual del precio\\[1pt]\arrayrulecolor{black!25}\midrule[0.25pt]
\end{tabularx}\endgroup\par
\end{minipage}\par\vspace{4pt}
\begin{minipage}{\linewidth}
\begingroup\fontsize{9}{10.8}\selectfont\begin{tabularx}{\linewidth}{@{}p{2.2cm}Y@{}}\arrayrulecolor{black}\toprule
\textbf{E-15-I1} & \textbf{Escenario inválido \textperiodcentered{} Prueba PA-15-I}\\[1pt]\midrule
\textbf{Justificación} & evita guardar un precio negativo.\\[1pt]\arrayrulecolor{black!25}\midrule[0.25pt]
\textbf{Análisis} & Probar también los límites de 0,001 y 20 euros por litro.\\[1pt]\arrayrulecolor{black!25}\midrule[0.25pt]
\textbf{Given} & Ana tiene Compacto con consumo de 6 l/100 km y precio manual de gasolina de 1,60 euros por litro\\[1pt]\arrayrulecolor{black!25}\midrule[0.25pt]
\textbf{When} & intenta fijar un precio de -1 euro por litro\\[1pt]\arrayrulecolor{black!25}\midrule[0.25pt]
\textbf{Then} & se rechaza el precio y el valor guardado permanece en 1,60 euros por litro\\[1pt]\arrayrulecolor{black!25}\midrule[0.25pt]
\end{tabularx}\endgroup\par
\end{minipage}\par\vspace{4pt}
\subsection*{HU-16. Estimar coste eléctrico}
\textbf{Requisito} R4\par
\textbf{Historia} Como usuario de un vehículo eléctrico quiero conocer el coste energético de una ruta para prever el gasto.\par
\begingroup\fontsize{8.5}{10.5}\selectfont\setlength{\tabcolsep}{4pt}
\begin{tabularx}{\linewidth}{@{}p{1.35cm}YY@{}}\toprule
\textbf{Escenario} & \textbf{Entrada o condición} & \textbf{Resultado}\\\midrule
E-16-V & Eléctrico, consumo y precio válidos & Coste según distancia\\[2pt]
E-16-I1 & Precio manual 0 euros/kWh & Precio rechazado\\[2pt]
E-16-V2 & Precio público no disponible & Ruta consultable sin coste\\[2pt]
\bottomrule\end{tabularx}\endgroup\par
\begin{minipage}{\linewidth}
\begingroup\fontsize{9}{10.8}\selectfont\begin{tabularx}{\linewidth}{@{}p{2.2cm}Y@{}}\arrayrulecolor{black}\toprule
\textbf{E-16-V} & \textbf{Escenario válido \textperiodcentered{} Prueba PA-16-V}\\[1pt]\midrule
\textbf{Justificación} & comprueba el cálculo en kWh y evita confundir las unidades del consumo y del precio.\\[1pt]\arrayrulecolor{black!25}\midrule[0.25pt]
\textbf{Análisis} & Revisar la conversión de euros por MWh a euros por kWh.\\[1pt]\arrayrulecolor{black!25}\midrule[0.25pt]
\textbf{Given} & Ana tiene Centro y UJI, Eléctrico con consumo de 15 kWh/100 km y precio manual de 0,20 \texteuro{}/kWh\\[1pt]\arrayrulecolor{black!25}\midrule[0.25pt]
\textbf{When} & calcula una ruta con Eléctrico entre Centro y UJI y el servicio de rutas está disponible\\[1pt]\arrayrulecolor{black!25}\midrule[0.25pt]
\textbf{Then} & el coste devuelto es (distancia en km / 100 \ensuremath{\times} 15 \ensuremath{\times} 0,20) \texteuro{} redondeado a céntimos, con el origen manual del precio\\[1pt]\arrayrulecolor{black!25}\midrule[0.25pt]
\end{tabularx}\endgroup\par
\end{minipage}\par\vspace{4pt}
\begin{minipage}{\linewidth}
\begingroup\fontsize{9}{10.8}\selectfont\begin{tabularx}{\linewidth}{@{}p{2.2cm}Y@{}}\arrayrulecolor{black}\toprule
\textbf{E-16-I1} & \textbf{Escenario inválido \textperiodcentered{} Prueba PA-16-I}\\[1pt]\midrule
\textbf{Justificación} & comprueba el mínimo positivo propuesto para el precio manual.\\[1pt]\arrayrulecolor{black!25}\midrule[0.25pt]
\textbf{Análisis} & Probar también un precio negativo y uno superior a 20 euros por kWh.\\[1pt]\arrayrulecolor{black!25}\midrule[0.25pt]
\textbf{Given} & Ana tiene Eléctrico y un precio manual de 0,20 euros por kWh\\[1pt]\arrayrulecolor{black!25}\midrule[0.25pt]
\textbf{When} & intenta guardar un precio de 0 euros por kWh\\[1pt]\arrayrulecolor{black!25}\midrule[0.25pt]
\textbf{Then} & se rechaza el precio y se conserva el valor de 0,20 euros por kWh\\[1pt]\arrayrulecolor{black!25}\midrule[0.25pt]
\end{tabularx}\endgroup\par
\end{minipage}\par\vspace{4pt}
\clearpage
\subsection*{HU-17. Estimar calorías}
\textbf{Requisito} R4\par
\textbf{Historia} Como usuario quiero estimar las calorías de una ruta a pie o en bicicleta para valorar el esfuerzo.\par
\begingroup\fontsize{8.5}{10.5}\selectfont\setlength{\tabcolsep}{4pt}
\begin{tabularx}{\linewidth}{@{}p{1.35cm}YY@{}}\toprule
\textbf{Escenario} & \textbf{Entrada o condición} & \textbf{Resultado}\\\midrule
E-17-V & A pie, peso 70 kg & Calorías estimadas\\[2pt]
E-17-I1 & Peso 0 kg & Peso rechazado\\[2pt]
E-17-V2 & Bicicleta, peso válido & Estimación con MET de bicicleta\\[2pt]
\bottomrule\end{tabularx}\endgroup\par
\begin{minipage}{\linewidth}
\begingroup\fontsize{9}{10.8}\selectfont\begin{tabularx}{\linewidth}{@{}p{2.2cm}Y@{}}\arrayrulecolor{black}\toprule
\textbf{E-17-V} & \textbf{Escenario válido \textperiodcentered{} Prueba PA-17-V}\\[1pt]\midrule
\textbf{Justificación} & comprueba las calorías a partir de la duración real obtenida y del peso configurado.\\[1pt]\arrayrulecolor{black!25}\midrule[0.25pt]
\textbf{Análisis} & Repetir en bicicleta y con los pesos extremos admitidos.\\[1pt]\arrayrulecolor{black!25}\midrule[0.25pt]
\textbf{Given} & Ana tiene Centro y UJI guardados, pesa 70 kg y el servicio de rutas está disponible\\[1pt]\arrayrulecolor{black!25}\midrule[0.25pt]
\textbf{When} & calcula una ruta a pie entre Centro y UJI\\[1pt]\arrayrulecolor{black!25}\midrule[0.25pt]
\textbf{Then} & se devuelven kcal estimadas con 3,5 \ensuremath{\times} 3,5 \ensuremath{\times} 70 / 200 \ensuremath{\times} duración en minutos, redondeadas al entero más próximo\\[1pt]\arrayrulecolor{black!25}\midrule[0.25pt]
\end{tabularx}\endgroup\par
\end{minipage}\par\vspace{4pt}
\begin{minipage}{\linewidth}
\begingroup\fontsize{9}{10.8}\selectfont\begin{tabularx}{\linewidth}{@{}p{2.2cm}Y@{}}\arrayrulecolor{black}\toprule
\textbf{E-17-I1} & \textbf{Escenario inválido \textperiodcentered{} Prueba PA-17-I}\\[1pt]\midrule
\textbf{Justificación} & un peso de cero no permite una estimación válida.\\[1pt]\arrayrulecolor{black!25}\midrule[0.25pt]
\textbf{Análisis} & Probar 19,99 y 300,01 kg y conservar el peso anterior.\\[1pt]\arrayrulecolor{black!25}\midrule[0.25pt]
\textbf{Given} & Ana tiene Centro y UJI guardados, peso configurado de 70 kg y el servicio de rutas disponible\\[1pt]\arrayrulecolor{black!25}\midrule[0.25pt]
\textbf{When} & intenta guardar un peso de 0 kg para calcular las calorías del trayecto\\[1pt]\arrayrulecolor{black!25}\midrule[0.25pt]
\textbf{Then} & se rechaza el valor y permanece el peso de 70 kg\\[1pt]\arrayrulecolor{black!25}\midrule[0.25pt]
\end{tabularx}\endgroup\par
\end{minipage}\par\vspace{4pt}
\subsection*{HU-18. Elegir el criterio de ruta}
\textbf{Requisito} R4\par
\textbf{Historia} Como usuario quiero priorizar duración, distancia o coste para escoger una alternativa adecuada.\par
\begingroup\fontsize{8.5}{10.5}\selectfont\setlength{\tabcolsep}{4pt}
\begin{tabularx}{\linewidth}{@{}p{1.35cm}YY@{}}\toprule
\textbf{Escenario} & \textbf{Entrada o condición} & \textbf{Resultado}\\\midrule
E-18-V & Más económica, alternativas y precio & Menor coste recibido primero\\[2pt]
E-18-I1 & Destino sin seleccionar & Cálculo rechazado\\[2pt]
E-18-V2 & Una sola alternativa disponible & Se muestra esa alternativa\\[2pt]
\bottomrule\end{tabularx}\endgroup\par
\begin{minipage}{\linewidth}
\begingroup\fontsize{9}{10.8}\selectfont\begin{tabularx}{\linewidth}{@{}p{2.2cm}Y@{}}\arrayrulecolor{black}\toprule
\textbf{E-18-V} & \textbf{Escenario válido \textperiodcentered{} Prueba PA-18-V}\\[1pt]\midrule
\textbf{Justificación} & compara los costes de las alternativas recibidas, sin suponer un recorrido fijo.\\[1pt]\arrayrulecolor{black!25}\midrule[0.25pt]
\textbf{Análisis} & Repetir con rapidez, distancia y alternativas empatadas.\\[1pt]\arrayrulecolor{black!25}\midrule[0.25pt]
\textbf{Given} & Ana tiene Centro, UJI y Compacto con consumo de 6 l/100 km, usa gasolina manual a 1,60 \texteuro{}/l y ORS está disponible\\[1pt]\arrayrulecolor{black!25}\midrule[0.25pt]
\textbf{When} & solicita una ruta en coche entre Centro y UJI con el criterio más económico\\[1pt]\arrayrulecolor{black!25}\midrule[0.25pt]
\textbf{Then} & la primera opción tiene coste menor o igual que cada alternativa obtenida y se muestra el origen del precio\\[1pt]\arrayrulecolor{black!25}\midrule[0.25pt]
\end{tabularx}\endgroup\par
\end{minipage}\par\vspace{4pt}
\begin{minipage}{\linewidth}
\begingroup\fontsize{9}{10.8}\selectfont\begin{tabularx}{\linewidth}{@{}p{2.2cm}Y@{}}\arrayrulecolor{black}\toprule
\textbf{E-18-I1} & \textbf{Escenario inválido \textperiodcentered{} Prueba PA-18-I}\\[1pt]\midrule
\textbf{Justificación} & sin destino no se puede comparar ninguna ruta.\\[1pt]\arrayrulecolor{black!25}\midrule[0.25pt]
\textbf{Análisis} & Comprobar también la ausencia de origen.\\[1pt]\arrayrulecolor{black!25}\midrule[0.25pt]
\textbf{Given} & Ana ha seleccionado Centro como origen y más económica como criterio, pero no ha seleccionado destino\\[1pt]\arrayrulecolor{black!25}\midrule[0.25pt]
\textbf{When} & solicita calcular la ruta\\[1pt]\arrayrulecolor{black!25}\midrule[0.25pt]
\textbf{Then} & se solicita seleccionar un destino y no se calcula una ruta nueva\\[1pt]\arrayrulecolor{black!25}\midrule[0.25pt]
\end{tabularx}\endgroup\par
\end{minipage}\par\vspace{4pt}
\clearpage
\subsection*{HU-19. Guardar una ruta}
\textbf{Requisito} R4\par
\textbf{Historia} Como usuario quiero guardar una ruta calculada para consultarla posteriormente.\par
\begingroup\fontsize{8.5}{10.5}\selectfont\setlength{\tabcolsep}{4pt}
\begin{tabularx}{\linewidth}{@{}p{1.35cm}YY@{}}\toprule
\textbf{Escenario} & \textbf{Entrada o condición} & \textbf{Resultado}\\\midrule
E-19-V & Resultado calculado y nombre válido & Ruta guardada\\[2pt]
E-19-I1 & Nombre de tres espacios & Ruta no guardada\\[2pt]
E-19-I2 & Sin resultado calculado & Guardado no disponible\\[2pt]
\bottomrule\end{tabularx}\endgroup\par
\begin{minipage}{\linewidth}
\begingroup\fontsize{9}{10.8}\selectfont\begin{tabularx}{\linewidth}{@{}p{2.2cm}Y@{}}\arrayrulecolor{black}\toprule
\textbf{E-19-V} & \textbf{Escenario válido \textperiodcentered{} Prueba PA-19-V}\\[1pt]\midrule
\textbf{Justificación} & comprueba que una ruta calculada puede guardarse para otro momento.\\[1pt]\arrayrulecolor{black!25}\midrule[0.25pt]
\textbf{Análisis} & Recuperar y comparar el trazado, las medidas y la fecha.\\[1pt]\arrayrulecolor{black!25}\midrule[0.25pt]
\textbf{Given} & Ana tiene un resultado real entre Centro y UJI calculado para su cuenta\\[1pt]\arrayrulecolor{black!25}\midrule[0.25pt]
\textbf{When} & lo guarda con el nombre Camino a clase\\[1pt]\arrayrulecolor{black!25}\midrule[0.25pt]
\textbf{Then} & se crea una ruta propia con ese nombre, geometría, métricas y fecha de cálculo, recuperable en su listado\\[1pt]\arrayrulecolor{black!25}\midrule[0.25pt]
\end{tabularx}\endgroup\par
\end{minipage}\par\vspace{4pt}
\begin{minipage}{\linewidth}
\begingroup\fontsize{9}{10.8}\selectfont\begin{tabularx}{\linewidth}{@{}p{2.2cm}Y@{}}\arrayrulecolor{black}\toprule
\textbf{E-19-I1} & \textbf{Escenario inválido \textperiodcentered{} Prueba PA-19-I}\\[1pt]\midrule
\textbf{Justificación} & una ruta guardada necesita un nombre no vacío.\\[1pt]\arrayrulecolor{black!25}\midrule[0.25pt]
\textbf{Análisis} & Probar nombre vacío y los límites de longitud.\\[1pt]\arrayrulecolor{black!25}\midrule[0.25pt]
\textbf{Given} & Ana tiene un resultado calculado entre Centro y UJI y ninguna ruta guardada\\[1pt]\arrayrulecolor{black!25}\midrule[0.25pt]
\textbf{When} & intenta guardarlo con un nombre de tres espacios\\[1pt]\arrayrulecolor{black!25}\midrule[0.25pt]
\textbf{Then} & se rechaza el nombre y la lista de rutas sigue vacía\\[1pt]\arrayrulecolor{black!25}\midrule[0.25pt]
\end{tabularx}\endgroup\par
\end{minipage}\par\vspace{4pt}
\subsection*{HU-20. Consultar rutas guardadas}
\textbf{Requisito} R4\par
\textbf{Historia} Como usuario quiero consultar mis rutas guardadas para recuperar un trayecto sin volver a calcularlo.\par
\begingroup\fontsize{8.5}{10.5}\selectfont\setlength{\tabcolsep}{4pt}
\begin{tabularx}{\linewidth}{@{}p{1.35cm}YY@{}}\toprule
\textbf{Escenario} & \textbf{Entrada o condición} & \textbf{Resultado}\\\midrule
E-20-V & Ruta guardada y sesión activa & Mismo trazado y métricas\\[2pt]
E-20-I1 & Sin sesión válida & Solicitar acceso\\[2pt]
E-20-V2 & Ninguna ruta guardada & Lista vacía\\[2pt]
\bottomrule\end{tabularx}\endgroup\par
\begin{minipage}{\linewidth}
\begingroup\fontsize{9}{10.8}\selectfont\begin{tabularx}{\linewidth}{@{}p{2.2cm}Y@{}}\arrayrulecolor{black}\toprule
\textbf{E-20-V} & \textbf{Escenario válido \textperiodcentered{} Prueba PA-20-V}\\[1pt]\midrule
\textbf{Justificación} & diferencia recuperar de recalcular con datos nuevos.\\[1pt]\arrayrulecolor{black!25}\midrule[0.25pt]
\textbf{Análisis} & Repetir sin conexión con el servicio de rutas.\\[1pt]\arrayrulecolor{black!25}\midrule[0.25pt]
\textbf{Given} & Ana tiene Camino a clase guardada con su geometría y sus métricas\\[1pt]\arrayrulecolor{black!25}\midrule[0.25pt]
\textbf{When} & abre su lista y selecciona Camino a clase\\[1pt]\arrayrulecolor{black!25}\midrule[0.25pt]
\textbf{Then} & recupera la misma geometría y las mismas métricas, identificadas como cálculo histórico\\[1pt]\arrayrulecolor{black!25}\midrule[0.25pt]
\end{tabularx}\endgroup\par
\end{minipage}\par\vspace{4pt}
\begin{minipage}{\linewidth}
\begingroup\fontsize{9}{10.8}\selectfont\begin{tabularx}{\linewidth}{@{}p{2.2cm}Y@{}}\arrayrulecolor{black}\toprule
\textbf{E-20-I1} & \textbf{Escenario inválido \textperiodcentered{} Prueba PA-20-I}\\[1pt]\midrule
\textbf{Justificación} & protege el historial de desplazamientos.\\[1pt]\arrayrulecolor{black!25}\midrule[0.25pt]
\textbf{Análisis} & Repetir con sesión caducada sin revelar ningún trayecto.\\[1pt]\arrayrulecolor{black!25}\midrule[0.25pt]
\textbf{Given} & no hay sesión válida\\[1pt]\arrayrulecolor{black!25}\midrule[0.25pt]
\textbf{When} & se solicita el listado de rutas guardadas\\[1pt]\arrayrulecolor{black!25}\midrule[0.25pt]
\textbf{Then} & se deniega el acceso sin revelar nombres ni trayectos\\[1pt]\arrayrulecolor{black!25}\midrule[0.25pt]
\end{tabularx}\endgroup\par
\end{minipage}\par\vspace{4pt}
\clearpage
\subsection*{HU-21. Eliminar una ruta guardada}
\textbf{Requisito} R4\par
\textbf{Historia} Como usuario quiero eliminar una ruta que no necesito para ordenar mi historial.\par
\begingroup\fontsize{8.5}{10.5}\selectfont\setlength{\tabcolsep}{4pt}
\begin{tabularx}{\linewidth}{@{}p{1.35cm}YY@{}}\toprule
\textbf{Escenario} & \textbf{Entrada o condición} & \textbf{Resultado}\\\midrule
E-21-V & Confirmar borrar Camino a clase & Solo queda Paseo\\[2pt]
E-21-I1 & Ruta eliminada antes de confirmar & Aviso, Paseo conservado\\[2pt]
E-21-V2 & Cancelar el borrado & Ruta conservada\\[2pt]
\bottomrule\end{tabularx}\endgroup\par
\begin{minipage}{\linewidth}
\begingroup\fontsize{9}{10.8}\selectfont\begin{tabularx}{\linewidth}{@{}p{2.2cm}Y@{}}\arrayrulecolor{black}\toprule
\textbf{E-21-V} & \textbf{Escenario válido \textperiodcentered{} Prueba PA-21-V}\\[1pt]\midrule
\textbf{Justificación} & delimita el efecto de la eliminación al historial.\\[1pt]\arrayrulecolor{black!25}\midrule[0.25pt]
\textbf{Análisis} & Repetir con ruta favorita y con cancelación del borrado.\\[1pt]\arrayrulecolor{black!25}\midrule[0.25pt]
\textbf{Given} & Ana tiene Camino a clase y Paseo\\[1pt]\arrayrulecolor{black!25}\midrule[0.25pt]
\textbf{When} & confirma eliminar Camino a clase\\[1pt]\arrayrulecolor{black!25}\midrule[0.25pt]
\textbf{Then} & solo permanece Paseo y no se alteran sus lugares ni vehículos\\[1pt]\arrayrulecolor{black!25}\midrule[0.25pt]
\end{tabularx}\endgroup\par
\end{minipage}\par\vspace{4pt}
\begin{minipage}{\linewidth}
\begingroup\fontsize{9}{10.8}\selectfont\begin{tabularx}{\linewidth}{@{}p{2.2cm}Y@{}}\arrayrulecolor{black}\toprule
\textbf{E-21-I1} & \textbf{Escenario inválido \textperiodcentered{} Prueba PA-21-I}\\[1pt]\midrule
\textbf{Justificación} & una ruta ya eliminada no puede borrarse de nuevo.\\[1pt]\arrayrulecolor{black!25}\midrule[0.25pt]
\textbf{Análisis} & Comprobar que las demás rutas se conservan.\\[1pt]\arrayrulecolor{black!25}\midrule[0.25pt]
\textbf{Given} & Ana tenía Camino a clase y Paseo, pero Camino a clase se ha eliminado desde otra sesión\\[1pt]\arrayrulecolor{black!25}\midrule[0.25pt]
\textbf{When} & confirma eliminar Camino a clase desde la lista anterior\\[1pt]\arrayrulecolor{black!25}\midrule[0.25pt]
\textbf{Then} & se informa de que ya no está disponible y Paseo permanece\\[1pt]\arrayrulecolor{black!25}\midrule[0.25pt]
\end{tabularx}\endgroup\par
\end{minipage}\par\vspace{4pt}
\subsection*{HU-22. Gestionar favoritos}
\textbf{Requisito} R5\par
\textbf{Historia} Como usuario quiero marcar y desmarcar favoritos para encontrar primero los elementos que uso con frecuencia.\par
\begingroup\fontsize{8.5}{10.5}\selectfont\setlength{\tabcolsep}{4pt}
\begin{tabularx}{\linewidth}{@{}p{1.35cm}YY@{}}\toprule
\textbf{Escenario} & \textbf{Entrada o condición} & \textbf{Resultado}\\\midrule
E-22-V & Marcar un elemento de cada lista & Favoritos primero, marcas guardadas\\[2pt]
E-22-I1 & Elemento eliminado antes de marcar & Aviso, resto conservado\\[2pt]
E-22-V2 & Desmarcar un favorito & Pierde su prioridad\\[2pt]
\bottomrule\end{tabularx}\endgroup\par
\begin{minipage}{\linewidth}
\begingroup\fontsize{9}{10.8}\selectfont\begin{tabularx}{\linewidth}{@{}p{2.2cm}Y@{}}\arrayrulecolor{black}\toprule
\textbf{E-22-V} & \textbf{Escenario válido \textperiodcentered{} Prueba PA-22-V}\\[1pt]\midrule
\textbf{Justificación} & comprueba la misma regla de favoritos en lugares, vehículos y rutas.\\[1pt]\arrayrulecolor{black!25}\midrule[0.25pt]
\textbf{Análisis} & Desmarcar cada favorito y volver a consultar las listas.\\[1pt]\arrayrulecolor{black!25}\midrule[0.25pt]
\textbf{Given} & Ana tiene sin marcar Centro y UJI, Compacto y Eléctrico, y Camino a clase y Paseo\\[1pt]\arrayrulecolor{black!25}\midrule[0.25pt]
\textbf{When} & marca UJI, Eléctrico y Paseo como favoritos y vuelve a consultar las tres listas\\[1pt]\arrayrulecolor{black!25}\midrule[0.25pt]
\textbf{Then} & cada elemento marcado conserva su estado y aparece antes que el no favorito de su lista\\[1pt]\arrayrulecolor{black!25}\midrule[0.25pt]
\end{tabularx}\endgroup\par
\end{minipage}\par\vspace{4pt}
\begin{minipage}{\linewidth}
\begingroup\fontsize{9}{10.8}\selectfont\begin{tabularx}{\linewidth}{@{}p{2.2cm}Y@{}}\arrayrulecolor{black}\toprule
\textbf{E-22-I1} & \textbf{Escenario inválido \textperiodcentered{} Prueba PA-22-I}\\[1pt]\midrule
\textbf{Justificación} & solo puede marcarse un elemento que todavía exista.\\[1pt]\arrayrulecolor{black!25}\midrule[0.25pt]
\textbf{Análisis} & Aplicar el mismo caso a vehículos y rutas.\\[1pt]\arrayrulecolor{black!25}\midrule[0.25pt]
\textbf{Given} & Ana tenía Centro y UJI, pero Centro se ha eliminado desde otra sesión\\[1pt]\arrayrulecolor{black!25}\midrule[0.25pt]
\textbf{When} & intenta marcar Centro como favorito desde la lista anterior\\[1pt]\arrayrulecolor{black!25}\midrule[0.25pt]
\textbf{Then} & se informa de que Centro ya no está disponible y UJI conserva su marca\\[1pt]\arrayrulecolor{black!25}\midrule[0.25pt]
\end{tabularx}\endgroup\par
\end{minipage}\par\vspace{4pt}
\clearpage
\subsection*{HU-23. Fijar transporte predeterminado}
\textbf{Requisito} R5\par
\textbf{Historia} Como usuario quiero fijar un modo y, si procede, un vehículo por defecto para preparar rutas con menos pasos.\par
\begingroup\fontsize{8.5}{10.5}\selectfont\setlength{\tabcolsep}{4pt}
\begin{tabularx}{\linewidth}{@{}p{1.35cm}YY@{}}\toprule
\textbf{Escenario} & \textbf{Entrada o condición} & \textbf{Resultado}\\\midrule
E-23-V & Coche y Compacto disponible & Predeterminados guardados\\[2pt]
E-23-I1 & Compacto eliminado antes de guardar & Referencia rechazada\\[2pt]
E-23-V2 & A pie o bicicleta & Sin vehículo asociado\\[2pt]
\bottomrule\end{tabularx}\endgroup\par
\begin{minipage}{\linewidth}
\begingroup\fontsize{9}{10.8}\selectfont\begin{tabularx}{\linewidth}{@{}p{2.2cm}Y@{}}\arrayrulecolor{black}\toprule
\textbf{E-23-V} & \textbf{Escenario válido \textperiodcentered{} Prueba PA-23-V}\\[1pt]\midrule
\textbf{Justificación} & reduce trabajo repetitivo sin bloquear la elección puntual.\\[1pt]\arrayrulecolor{black!25}\midrule[0.25pt]
\textbf{Análisis} & Repetir con bicicleta y a pie, sin vehículo asociado.\\[1pt]\arrayrulecolor{black!25}\midrule[0.25pt]
\textbf{Given} & Compacto pertenece a Ana\\[1pt]\arrayrulecolor{black!25}\midrule[0.25pt]
\textbf{When} & guarda coche y Compacto como valores predeterminados y abre una planificación nueva\\[1pt]\arrayrulecolor{black!25}\midrule[0.25pt]
\textbf{Then} & se seleccionan coche y Compacto, permitiendo cambiarlos para esa consulta\\[1pt]\arrayrulecolor{black!25}\midrule[0.25pt]
\end{tabularx}\endgroup\par
\end{minipage}\par\vspace{4pt}
\begin{minipage}{\linewidth}
\begingroup\fontsize{9}{10.8}\selectfont\begin{tabularx}{\linewidth}{@{}p{2.2cm}Y@{}}\arrayrulecolor{black}\toprule
\textbf{E-23-I1} & \textbf{Escenario inválido \textperiodcentered{} Prueba PA-23-I}\\[1pt]\midrule
\textbf{Justificación} & no se puede guardar una referencia a un vehículo eliminado.\\[1pt]\arrayrulecolor{black!25}\midrule[0.25pt]
\textbf{Análisis} & Comprobar que el modo anterior se conserva.\\[1pt]\arrayrulecolor{black!25}\midrule[0.25pt]
\textbf{Given} & Ana usa a pie por defecto y tiene Compacto seleccionado en la configuración, pero el vehículo se elimina desde otra sesión antes de guardar\\[1pt]\arrayrulecolor{black!25}\midrule[0.25pt]
\textbf{When} & intenta guardar coche y Compacto como predeterminados\\[1pt]\arrayrulecolor{black!25}\midrule[0.25pt]
\textbf{Then} & se informa de que Compacto ya no está disponible y se conserva a pie como modo predeterminado\\[1pt]\arrayrulecolor{black!25}\midrule[0.25pt]
\end{tabularx}\endgroup\par
\end{minipage}\par\vspace{4pt}
\subsection*{HU-24. Fijar criterio predeterminado}
\textbf{Requisito} R5\par
\textbf{Historia} Como usuario quiero guardar un criterio de ruta por defecto para no repetir mi preferencia.\par
\begingroup\fontsize{8.5}{10.5}\selectfont\setlength{\tabcolsep}{4pt}
\begin{tabularx}{\linewidth}{@{}p{1.35cm}YY@{}}\toprule
\textbf{Escenario} & \textbf{Entrada o condición} & \textbf{Resultado}\\\midrule
E-24-V & Guardar más corta & Criterio inicial más corta\\[2pt]
E-24-I1 & Almacenamiento no disponible & Error, preferencia conservada\\[2pt]
E-24-V2 & Cambiar criterio solo para una consulta & Predeterminado conservado\\[2pt]
\bottomrule\end{tabularx}\endgroup\par
\begin{minipage}{\linewidth}
\begingroup\fontsize{9}{10.8}\selectfont\begin{tabularx}{\linewidth}{@{}p{2.2cm}Y@{}}\arrayrulecolor{black}\toprule
\textbf{E-24-V} & \textbf{Escenario válido \textperiodcentered{} Prueba PA-24-V}\\[1pt]\midrule
\textbf{Justificación} & comprueba el efecto posterior de guardar, además de la notificación.\\[1pt]\arrayrulecolor{black!25}\midrule[0.25pt]
\textbf{Análisis} & Repetir con las tres opciones de criterio admitidas.\\[1pt]\arrayrulecolor{black!25}\midrule[0.25pt]
\textbf{Given} & Ana tiene más rápida como criterio predeterminado\\[1pt]\arrayrulecolor{black!25}\midrule[0.25pt]
\textbf{When} & guarda más corta y abre una planificación nueva\\[1pt]\arrayrulecolor{black!25}\midrule[0.25pt]
\textbf{Then} & aparece más corta como criterio inicial y la preferencia se conserva\\[1pt]\arrayrulecolor{black!25}\midrule[0.25pt]
\end{tabularx}\endgroup\par
\end{minipage}\par\vspace{4pt}
\begin{minipage}{\linewidth}
\begingroup\fontsize{9}{10.8}\selectfont\begin{tabularx}{\linewidth}{@{}p{2.2cm}Y@{}}\arrayrulecolor{black}\toprule
\textbf{E-24-I1} & \textbf{Escenario inválido \textperiodcentered{} Prueba PA-24-I}\\[1pt]\midrule
\textbf{Justificación} & el fallo de guardado no debe mostrarse como un cambio confirmado.\\[1pt]\arrayrulecolor{black!25}\midrule[0.25pt]
\textbf{Análisis} & Volver a abrir y comprobar la preferencia anterior.\\[1pt]\arrayrulecolor{black!25}\midrule[0.25pt]
\textbf{Given} & Ana tiene más corta como criterio predeterminado y el almacenamiento deja de estar disponible\\[1pt]\arrayrulecolor{black!25}\midrule[0.25pt]
\textbf{When} & intenta guardar más rápida\\[1pt]\arrayrulecolor{black!25}\midrule[0.25pt]
\textbf{Then} & se informa del fallo de guardado y al recuperar el almacenamiento sigue guardado más corta\\[1pt]\arrayrulecolor{black!25}\midrule[0.25pt]
\end{tabularx}\endgroup\par
\end{minipage}\par\vspace{4pt}
\clearpage
\subsection*{HU-25. Recuperar estado tras un cierre}
\textbf{Requisito} R6\par
\textbf{Historia} Como usuario quiero recuperar contenidos y aspecto al volver a abrir la aplicación para evitar reconfigurarla.\par
\begingroup\fontsize{8.5}{10.5}\selectfont\setlength{\tabcolsep}{4pt}
\begin{tabularx}{\linewidth}{@{}p{1.35cm}YY@{}}\toprule
\textbf{Escenario} & \textbf{Entrada o condición} & \textbf{Resultado}\\\midrule
E-25-V & Cierre normal con datos guardados & Estado recuperado\\[2pt]
E-25-I1 & Almacenamiento inaccesible & Error de carga\\[2pt]
E-25-V2 & Listas vacías guardadas & Estado vacío recuperado\\[2pt]
\bottomrule\end{tabularx}\endgroup\par
\begin{minipage}{\linewidth}
\begingroup\fontsize{9}{10.8}\selectfont\begin{tabularx}{\linewidth}{@{}p{2.2cm}Y@{}}\arrayrulecolor{black}\toprule
\textbf{E-25-V} & \textbf{Escenario válido \textperiodcentered{} Prueba PA-25-V}\\[1pt]\midrule
\textbf{Justificación} & comprueba la conservación del estado confirmado al cerrar y volver a abrir la aplicación.\\[1pt]\arrayrulecolor{black!25}\midrule[0.25pt]
\textbf{Análisis} & Comparar datos, favoritos y preferencias guardados.\\[1pt]\arrayrulecolor{black!25}\midrule[0.25pt]
\textbf{Given} & Ana ha guardado Centro, UJI, Compacto y Camino a clase, ha marcado UJI como favorito y usa tema oscuro y el último trayecto Centro-UJI a pie con criterio más corta\\[1pt]\arrayrulecolor{black!25}\midrule[0.25pt]
\textbf{When} & cierra la aplicación y vuelve a abrirla sobre la misma base de datos, iniciando sesión si es necesario\\[1pt]\arrayrulecolor{black!25}\midrule[0.25pt]
\textbf{Then} & recupera los mismos registros, favoritos, tema y selección del último trayecto\\[1pt]\arrayrulecolor{black!25}\midrule[0.25pt]
\end{tabularx}\endgroup\par
\end{minipage}\par\vspace{4pt}
\begin{minipage}{\linewidth}
\begingroup\fontsize{9}{10.8}\selectfont\begin{tabularx}{\linewidth}{@{}p{2.2cm}Y@{}}\arrayrulecolor{black}\toprule
\textbf{E-25-I1} & \textbf{Escenario inválido \textperiodcentered{} Prueba PA-25-I}\\[1pt]\midrule
\textbf{Justificación} & un fallo de almacenamiento debe comunicarse sin mostrar datos vacíos como si fueran reales.\\[1pt]\arrayrulecolor{black!25}\midrule[0.25pt]
\textbf{Análisis} & Repetir con base inaccesible al arrancar el servidor.\\[1pt]\arrayrulecolor{black!25}\midrule[0.25pt]
\textbf{Given} & Ana tiene Centro, Compacto y Camino a clase guardados, pero el almacenamiento no se puede abrir\\[1pt]\arrayrulecolor{black!25}\midrule[0.25pt]
\textbf{When} & intenta iniciar la aplicación\\[1pt]\arrayrulecolor{black!25}\midrule[0.25pt]
\textbf{Then} & se informa del fallo de almacenamiento y no se presenta una lista vacía como si fuera su estado real\\[1pt]\arrayrulecolor{black!25}\midrule[0.25pt]
\end{tabularx}\endgroup\par
\end{minipage}\par\vspace{4pt}
\subsection*{HU-26. Recuperar estado tras interrupción}
\textbf{Requisito} R6\par
\textbf{Historia} Como usuario quiero conservar los datos confirmados tras un cierre involuntario para no perder el trabajo guardado.\par
\begingroup\fontsize{8.5}{10.5}\selectfont\setlength{\tabcolsep}{4pt}
\begin{tabularx}{\linewidth}{@{}p{1.35cm}YY@{}}\toprule
\textbf{Escenario} & \textbf{Entrada o condición} & \textbf{Resultado}\\\midrule
E-26-V & Interrupción tras confirmar Centro & Centro conservado\\[2pt]
E-26-I1 & Interrupción antes de confirmar UJI & Centro conservado, UJI ausente\\[2pt]
E-26-V2 & Interrupción tras confirmar ambos & Ambos conservados\\[2pt]
\bottomrule\end{tabularx}\endgroup\par
\begin{minipage}{\linewidth}
\begingroup\fontsize{9}{10.8}\selectfont\begin{tabularx}{\linewidth}{@{}p{2.2cm}Y@{}}\arrayrulecolor{black}\toprule
\textbf{E-26-V} & \textbf{Escenario válido \textperiodcentered{} Prueba PA-26-V}\\[1pt]\midrule
\textbf{Justificación} & comprueba que una interrupción no borra un dato ya guardado.\\[1pt]\arrayrulecolor{black!25}\midrule[0.25pt]
\textbf{Análisis} & Reabrir la base en otro proceso y comprobar duplicados.\\[1pt]\arrayrulecolor{black!25}\midrule[0.25pt]
\textbf{Given} & el servidor ha confirmado el alta de Centro en una base de datos en disco\\[1pt]\arrayrulecolor{black!25}\midrule[0.25pt]
\textbf{When} & se fuerza el cierre del proceso con SIGKILL y se abre otro sobre la misma base\\[1pt]\arrayrulecolor{black!25}\midrule[0.25pt]
\textbf{Then} & Centro permanece íntegro y se puede consultar sin duplicados\\[1pt]\arrayrulecolor{black!25}\midrule[0.25pt]
\end{tabularx}\endgroup\par
\end{minipage}\par\vspace{4pt}
\begin{minipage}{\linewidth}
\begingroup\fontsize{9}{10.8}\selectfont\begin{tabularx}{\linewidth}{@{}p{2.2cm}Y@{}}\arrayrulecolor{black}\toprule
\textbf{E-26-I1} & \textbf{Escenario inválido \textperiodcentered{} Prueba PA-26-I}\\[1pt]\midrule
\textbf{Justificación} & el guardado interrumpido no puede considerarse terminado ni dejar datos parciales.\\[1pt]\arrayrulecolor{black!25}\midrule[0.25pt]
\textbf{Análisis} & Interrumpir antes de confirmar y comprobar los datos anteriores.\\[1pt]\arrayrulecolor{black!25}\midrule[0.25pt]
\textbf{Given} & Centro está confirmado y una transacción ha escrito UJI, pero aún no ha ejecutado COMMIT\\[1pt]\arrayrulecolor{black!25}\midrule[0.25pt]
\textbf{When} & el proceso se interrumpe antes del COMMIT y después se abre la misma base de datos\\[1pt]\arrayrulecolor{black!25}\midrule[0.25pt]
\textbf{Then} & Centro se conserva y UJI no aparece como alta confirmada. No hay datos parciales\\[1pt]\arrayrulecolor{black!25}\midrule[0.25pt]
\end{tabularx}\endgroup\par
\end{minipage}\par\vspace{4pt}
\endgroup
